\documentclass[11pt,a4paper]{article}

% --- Page margins and paragraphs ---
\usepackage[margin=1in]{geometry}
\usepackage{parskip}

% --- Basic math and typesetting ---
\usepackage{amsmath,amssymb,amsthm}
\usepackage{xcolor} % 

% --- Theorem and definition environments ---
\theoremstyle{definition}
\newtheorem{definition}{Definition}[section]
\newtheorem{example}{Example}[section]

\theoremstyle{plain}
\newtheorem{theorem}{Theorem}[section]
\newtheorem{proposition}{Proposition}[section]
\newtheorem{lemma}{Lemma}[section]

\begin{document}

% ==================================================
% Title and author information
% ==================================================
\begin{center}
    {\LARGE \textbf{Mathematics Method for Physics--Lecture 17}} \\[0.6em]
    {\large The Inverse Laplace Transform and Its Applications}
\end{center}

\vspace{-0.5em}
\noindent\rule{\textwidth}{0.5pt}
\vspace{0.3em}

\begin{center}
    Instructor: Zhaolong Han, SOCE, Shanghai Jiao Tong University;
\end{center}

\vspace{-0.6em}
\noindent Email: han.arkey@sjtu.edu.cn

\vspace{0.4em}
\noindent\rule{\textwidth}{0.5pt}

\vspace{0.8em}
\noindent\textit{Conventions: $p=\sigma+\mathrm{i}\omega$; $u(t)$ is the unit step function;
$\mathrm{e},\mathrm{i},\mathrm{d}$ upright; formulas are numbered (17.1)--(17.5).}

\vspace{1.5em}

\setcounter{section}{16}

% ==================================================
% Main content
% ==================================================
\section{Inverse Laplace Transform and Applications}

The core of integral-transform theory can be summarized in four items:
\begin{enumerate}\setlength{\itemsep}{2pt}\setlength{\topsep}{3pt}
    \item \textbf{Definition and theorems}: the Fourier transform (on $(-\infty,+\infty)$) and the
          Laplace transform (on $[0,+\infty)$), together with their existence theorems;
    \item \textbf{Common transform pairs}: elementary functions;
    \item \textbf{Properties}: linearity, similarity, delay (time shift), frequency shift,
          differentiation, integration, periodicity --- about seven or eight basic properties
          (here too ``seven or eight'' is approximate; the definitive list is the one collected in
          Chapter 16);
    \item \textbf{Applications}: solving differential and integral equations, and evaluating
          certain generalized integrals.
\end{enumerate}

\subsection{The Inverse Laplace Transform}

\subsubsection{The complex inversion integral formula}

Everything rests on one observation: the Laplace transform of $f(t)$ is exactly the Fourier
transform of $f(t)u(t)\,\mathrm{e}^{-\sigma t}$, where $\sigma$ is a real constant. Indeed, $u(t)$
pulls the lower limit from $-\infty$ to $0$,
\begin{equation*}
\mathcal{F}\left[f(t)u(t)\,\mathrm{e}^{-\sigma t}\right]
=\int_{-\infty}^{+\infty}f(t)u(t)\,\mathrm{e}^{-\sigma t}\,\mathrm{e}^{-\mathrm{i}\omega t}\,\mathrm{d}t
=\int_{0}^{+\infty}f(t)\,\mathrm{e}^{-(\sigma+\mathrm{i}\omega)t}\,\mathrm{d}t=F(p),
\qquad p=\sigma+\mathrm{i}\omega .
\end{equation*}

Applying the Fourier integral formula to $f(t)u(t)\,\mathrm{e}^{-\sigma t}$ gives
\begin{equation*}
f(t)u(t)\,\mathrm{e}^{-\sigma t}=\frac{1}{2\pi}\int_{-\infty}^{+\infty}
F(\sigma+\mathrm{i}\omega)\,\mathrm{e}^{\mathrm{i}\omega t}\,\mathrm{d}\omega .
\end{equation*}

Multiply by $\mathrm{e}^{\sigma t}$ and use $u(t)=1$ for $t>0$; then substitute
$p=\sigma+\mathrm{i}\omega$, so that $\mathrm{d}p=\mathrm{i}\,\mathrm{d}\omega$ and
$\omega:-\infty\to+\infty$ corresponds to
$p:\sigma-\mathrm{i}\infty\to\sigma+\mathrm{i}\infty$. This yields
\begin{equation}
f(t)=\frac{1}{2\pi}\int_{-\infty}^{+\infty}F(\sigma+\mathrm{i}\omega)\,
\mathrm{e}^{(\sigma+\mathrm{i}\omega)t}\,\mathrm{d}\omega
=\frac{1}{2\pi\mathrm{i}}\int_{\sigma-\mathrm{i}\infty}^{\sigma+\mathrm{i}\infty}
F(p)\,\mathrm{e}^{pt}\,\mathrm{d}p\qquad(t>0). \tag{17.1}
\end{equation}

The factor in front differs between the two forms ($1/2\pi$ versus $1/(2\pi\mathrm{i})$); the
difference is precisely $\mathrm{d}p=\mathrm{i}\,\mathrm{d}\omega$. The integral on the right is
called the \textbf{Bromwich integral}; its path is any line $\operatorname{Re}p=\sigma$ parallel to
the imaginary axis, lying in the region of convergence.

\begin{theorem}[Complex inversion integral]
Let $f(t)$ be piecewise smooth on $t\ge 0$ and of exponential order,
$\left|f(t)\right|\le M\mathrm{e}^{\sigma_c t}$, and let
$\mathcal{L}\left[f(t)\right]=F(p)$. Take $\sigma>\sigma_c$. Then for $t>0$
\begin{equation}
f(t)=\frac{1}{2\pi\mathrm{i}}\int_{\sigma-\mathrm{i}\infty}^{\sigma+\mathrm{i}\infty}
F(p)\,\mathrm{e}^{pt}\,\mathrm{d}p , \tag{17.2}
\end{equation}
where at a discontinuity of $f$ the left-hand side is to be replaced by the mean of the one-sided
limits, $\frac{f(t+0)+f(t-0)}{2}$.
\end{theorem}

The condition $\sigma>\sigma_c$ makes $f(t)u(t)\,\mathrm{e}^{-\sigma t}$ absolutely integrable on
$(-\infty,+\infty)$, which is what the Fourier integral formula requires; it also puts the path
inside the region of convergence $\operatorname{Re}p>\sigma_0$ ($\sigma_0$ is the abscissa of
convergence, $\sigma_0\le\sigma_c$). The image function must in addition satisfy
$\lim_{p\to\infty}F(p)=0$ --- the necessary condition for a Laplace transform (Chapter 16,
remark (1)).

\subsubsection{Inversion by the residue theorem}

The complex integral (17.2) is rarely evaluated directly; instead the path is closed to the left and
the integral becomes a sum of residues.

Close the path with the left semicircle $C_{\rho}$ of radius $\rho$ centred at
$\sigma+\mathrm{i}0$, which together with the vertical segment bounds a left half-disc. On
$C_{\rho}$ write $p=\sigma+\rho\,\mathrm{e}^{\mathrm{i}\theta}$ with
$\pi/2<\theta<3\pi/2$; then
\begin{equation*}
\left|\mathrm{e}^{pt}\right|=\mathrm{e}^{\sigma t}\mathrm{e}^{\rho t\cos\theta},
\qquad \cos\theta<0 .
\end{equation*}

The exponential factor decays along the arc, and the integral over the arc is bounded by
$\pi\mathrm{e}^{\sigma t}/t$ (the mechanism behind Jordan's lemma). Hence, provided
$F(p)$ tends to zero uniformly on the arc --- guaranteed by $\lim_{p\to\infty}F(p)=0$ --- the arc
integral tends to zero as $\rho\to+\infty$. For $t<0$ the exponential factor forces the path to be
closed to the right, where there are no singularities, so the integral vanishes; this is why (17.2)
is stated for $t>0$ only.

\begin{theorem}[Inversion by residues]
Let $F(p)$ have finitely many singular points $p_1,\dots,p_n$, all to the left of the line
$\operatorname{Re}p=\sigma$, and let $\lim_{p\to\infty}F(p)=0$. Then for $t>0$
\begin{equation}
f(t)=\mathcal{L}^{-1}\left[F(p)\right]
=\sum_{k=1}^{n}\operatorname{Res}\left[F(p)\,\mathrm{e}^{pt},\,p_k\right]. \tag{17.3}
\end{equation}
\end{theorem}

Proof. Combine the path of (17.2) with the left semicircle into a counterclockwise closed
contour containing every singularity of $F(p)\mathrm{e}^{pt}$. By the residue theorem the contour
integral equals $2\pi\mathrm{i}\sum_{k}\operatorname{Res}$; letting $\rho\to+\infty$, the arc
integral vanishes and the vertical segment tends to the path of (17.2), which gives (17.3).
\hfill$\blacksquare$

For a rational function $F(p)=\dfrac{A(p)}{B(p)}$ with $\deg B>\deg A$ all singularities are poles
and $F(\infty)=0$ holds automatically, so this is the case met in practice. If every $p_k$ is a
simple pole,
\begin{equation}
f(t)=\sum_{k=1}^{n}\frac{A(p_k)}{B'(p_k)}\,\mathrm{e}^{p_k t} ; \tag{17.4}
\end{equation}
if $p_k$ is a pole of order $m_k$,
\begin{equation}
f(t)=\sum_{k=1}^{n}\frac{1}{(m_k-1)!}\lim_{p\to p_k}\frac{\mathrm{d}^{m_k-1}}
{\mathrm{d}p^{m_k-1}}\left[(p-p_k)^{m_k}\frac{A(p)}{B(p)}\,\mathrm{e}^{pt}\right]. \tag{17.5}
\end{equation}

\begin{samepage}
\begin{example}
Find $\mathcal{L}^{-1}\left[\dfrac{1}{p^2+1}\right]$.
\end{example}

Solution. The poles $p=\pm\mathrm{i}$ are simple; by (17.4),
\begin{equation*}
\mathcal{L}^{-1}\left[\frac{1}{p^2+1}\right]
=\frac{\mathrm{e}^{\mathrm{i}t}}{2\mathrm{i}}+\frac{\mathrm{e}^{-\mathrm{i}t}}{2(-\mathrm{i})}
=\sin t .
\end{equation*}
\end{samepage}

\begin{example}
Solve the initial-value problem
\begin{equation*}
x''(t)-2x'(t)+2x(t)=2\mathrm{e}^{t}\cos t,\qquad x(0)=0,\quad x'(0)=0 .
\end{equation*}
\end{example}

Solution. Transforming ($\mathcal{L}[x'']=p^2X$, $\mathcal{L}[x']=pX$,
$\mathcal{L}\left[2\mathrm{e}^{t}\cos t\right]=\dfrac{2(p-1)}{(p-1)^2+1}$) gives
\begin{equation*}
X(p)=\frac{2(p-1)}{\left[(p-1)^2+1\right]^2} .
\end{equation*}

\textbf{Method 1 (residues).} The points $p=1\pm\mathrm{i}$ are poles of \textbf{order two}, so
(17.5) applies. For $p_0=1+\mathrm{i}$,
\begin{equation*}
\operatorname{Res}\left[X(p)\mathrm{e}^{pt},1+\mathrm{i}\right]
=\lim_{p\to 1+\mathrm{i}}\frac{\mathrm{d}}{\mathrm{d}p}\left[
\frac{2(p-1)\mathrm{e}^{pt}}{\left(p-1+\mathrm{i}\right)^2}\right],
\end{equation*}
and after differentiating, substituting $p=1+\mathrm{i}$ and adding the residue at
$p=1-\mathrm{i}$, the answer is again $\mathrm{e}^{t}t\sin t$. The differentiation and substitution
for a second-order pole are laborious --- that is the drawback of this route.

\textbf{Method 2 (properties).} Put $q=p-1$; what has to be inverted is
$\mathcal{L}^{-1}\left[\frac{2q}{\left(q^2+1\right)^2}\right]$. By the derivative-of-the-image
property, $\dfrac{\mathrm{d}}{\mathrm{d}p}\dfrac{1}{p^2+1}
=-\dfrac{2p}{\left(p^2+1\right)^2}$, i.e.
\begin{equation*}
\mathcal{L}\left[t\sin t\right]=-\left(\frac{1}{p^2+1}\right)'
=\frac{2p}{\left(p^2+1\right)^2},
\end{equation*}
so $\mathcal{L}^{-1}\left[\frac{2q}{\left(q^2+1\right)^2}\right]=t\sin t$. Then the shift property
$\mathcal{L}^{-1}\left[F(p-p_0)\right]=\mathrm{e}^{p_0 t}\mathcal{L}^{-1}\left[F(p)\right]$ with
$p_0=1$ restores $q\to p-1$:
\begin{equation*}
x(t)=\mathrm{e}^{t}\,t\sin t .
\end{equation*}

Both methods agree; whenever the properties suffice, the residues are not worth the work.

\clearpage
\setlength{\parskip}{4pt plus 1pt}
\setlength{\abovedisplayskip}{5pt plus 1pt minus 1pt}
\setlength{\belowdisplayskip}{5pt plus 1pt minus 1pt}
\subsection{Applications of the Laplace Transform}

Solving a linear differential (integral) equation by the Laplace transform takes three steps:
\begin{enumerate}\setlength{\itemsep}{2pt}\setlength{\topsep}{3pt}
    \item transform the equation, turning the differential (integral) equation into an
          \textbf{algebraic} equation for the image function;
    \item solve that algebraic equation for $X(p)$;
    \item invert $X(p)$ to obtain the solution.
\end{enumerate}

Step 3 usually splits $X(p)$ into table entries by partial fractions, or uses the residue formula
of \S 17.1.2 (17.3).

\subsubsection{Linear differential (and integral) equations}

\begin{example}
Solve the differential-integral equation
\begin{equation*}
x'(t)-2\int_{0}^{t}u(s)\,x(t-s)\,\mathrm{d}s+3\int_{0}^{t}x(s)\,\mathrm{d}s=u(t-1),
\qquad x(0)=0 ,
\end{equation*}
where $u(t)$ is the unit step function.
\end{example}

Solution. Transform the four terms separately:
\begin{itemize}\setlength{\itemsep}{2pt}\setlength{\topsep}{3pt}
    \item first term: $\mathcal{L}\left[x'\right]=pX(p)-x(0)=pX$;
    \item second term: $\displaystyle\int_{0}^{t}u(s)x(t-s)\,\mathrm{d}s=u(t)*x(t)$ is a
          convolution, so by the convolution theorem
          $\mathcal{L}\left[u*x\right]=\dfrac{X(p)}{p}$ (since $\mathcal{L}[u]=\dfrac{1}{p}$);
    \item third term: by the integration property
          $\mathcal{L}\left[\displaystyle\int_{0}^{t}x(s)\,\mathrm{d}s\right]=\dfrac{X(p)}{p}$;
    \item right-hand side: $\mathcal{L}\left[u(t-1)\right]=\dfrac{\mathrm{e}^{-p}}{p}$.
\end{itemize}

Hence
\begin{equation*}
pX-\frac{2}{p}X+\frac{3}{p}X=\frac{\mathrm{e}^{-p}}{p},
\end{equation*}
and multiplying by $p$ gives $\left(p^2+1\right)X=\mathrm{e}^{-p}$, i.e.
\begin{equation*}
X(p)=\frac{\mathrm{e}^{-p}}{p^2+1} .
\end{equation*}

By the delay property the factor $\mathrm{e}^{-p}$ simply shifts $\sin t$ to the right by 1:
\begin{equation*}
x(t)=u(t-1)\sin(t-1).
\end{equation*}

The factor $u(t-1)$ cannot be omitted: for $t<1$ both the right-hand side and the initial data
vanish, so the solution must vanish too, while for $t\ge 1$ it is a sine oscillation of unit
amplitude.

The same three steps cover systems of constant-coefficient equations (reduced to a linear algebraic
system for the images), boundary-value problems with a condition at one end only (the unknown
initial value at the other end is carried as a parameter and fixed afterwards), and circuit
equations such as the series $RLC$ circuit (where the inductance and capacitance contribute $LpI$
and $I/(Cp)$).

\subsubsection{Generalized integrals}

The integration and differentiation properties of the image function yield two classes of
generalized integrals.

\textbf{Type 1 (containing $1/t$).} Let $F(p)=\mathcal{L}[f(t)]$. The integration property of the
image function gives $\mathcal{L}\left[\dfrac{f(t)}{t}\right]
=\displaystyle\int_{p}^{+\infty}F(s)\,\mathrm{d}s$. Letting $p\to 0^{+}$ along the real axis inside
the region of convergence, and using the continuity of the Laplace integral at a point of
convergence (Abel's theorem),
\begin{equation*}
\int_{0}^{+\infty}\frac{f(t)}{t}\,\mathrm{d}t
=\lim_{p\to 0^{+}}\int_{p}^{+\infty}F(s)\,\mathrm{d}s
=\int_{0}^{+\infty}F(s)\,\mathrm{d}s .
\end{equation*}

Setting $\mathrm{e}^{-pt}=1$ in the integrand is not a licence taken lightly: its legitimacy is
exactly Abel's theorem.

\begin{example}
Evaluate the Dirichlet integral
$\displaystyle\int_{0}^{+\infty}\frac{\sin t}{t}\,\mathrm{d}t$.
\end{example}

Solution. Take $f(t)=\sin t$, $F(p)=\dfrac{1}{p^2+1}$; by Type 1,
\begin{equation*}
\int_{0}^{+\infty}\frac{\sin t}{t}\,\mathrm{d}t
=\int_{0}^{+\infty}\frac{\mathrm{d}s}{1+s^2}
=\left.\arctan s\right|_{0}^{+\infty}=\frac{\pi}{2} .
\end{equation*}

\textbf{Type 2 (containing $t^{n}$).} The derivative-of-the-image property gives
\begin{equation*}
F^{(n)}(p)=(-1)^{n}\int_{0}^{+\infty}t^{n}f(t)\,\mathrm{e}^{-pt}\,\mathrm{d}t ,
\quad\Longrightarrow\quad
\int_{0}^{+\infty}t^{n}f(t)\,\mathrm{d}t=(-1)^{n}F^{(n)}(0)\quad(p\to 0),
\end{equation*}
where the step to the right needs $F$ analytic on $\operatorname{Re}p\ge 0$ and the integral
convergent.

\begin{example}
Evaluate
$\displaystyle\int_{0}^{+\infty}t^{2}\mathrm{e}^{-\alpha t}\sin\beta t\,\mathrm{d}t
\quad(\alpha,\beta>0)$.
\end{example}

Solution. Take $f(t)=\mathrm{e}^{-\alpha t}\sin\beta t$,
$F(p)=\dfrac{\beta}{(p+\alpha)^2+\beta^2}$; by Type 2 with $n=2$,
\begin{equation*}
\int_{0}^{+\infty}t^{2}\mathrm{e}^{-\alpha t}\sin\beta t\,\mathrm{d}t=F''(0)
=\left.\frac{-2\beta\left[\beta^2-3(p+\alpha)^2\right]}
{\left[(p+\alpha)^2+\beta^2\right]^3}\right|_{p=0}
=\frac{2\beta\left(3\alpha^2-\beta^2\right)}{\left(\alpha^2+\beta^2\right)^3} .
\end{equation*}

For $\beta>\sqrt{3}\,\alpha$ the result is negative: the rapid oscillation of $\sin\beta t$ makes
the positive and negative parts almost cancel. This is not a computational slip.

\textbf{Complex-parameter method.} Writing $\sin\beta t=\operatorname{Im}\mathrm{e}^{\mathrm{i}\beta t}$,
the integral is the image function evaluated at $p=\alpha-\mathrm{i}\beta$:
\begin{equation*}
\int_{0}^{+\infty}t^{2}\mathrm{e}^{-\alpha t}\mathrm{e}^{\mathrm{i}\beta t}\,\mathrm{d}t
=\left.\frac{2}{p^{3}}\right|_{p=\alpha-\mathrm{i}\beta}
=\frac{2}{(\alpha-\mathrm{i}\beta)^{3}},
\end{equation*}
whose imaginary part reproduces the result above. In general, for real-valued $f$ and $\alpha$
inside the region of convergence,
\begin{equation*}
\int_{0}^{+\infty}f(t)\,\mathrm{e}^{-\alpha t}\sin\beta t\,\mathrm{d}t
=\operatorname{Im}F(\alpha-\mathrm{i}\beta),\qquad
\int_{0}^{+\infty}f(t)\,\mathrm{e}^{-\alpha t}\cos\beta t\,\mathrm{d}t
=\operatorname{Re}F(\alpha-\mathrm{i}\beta).
\end{equation*}

This is the standard way to evaluate generalized integrals carrying a damping factor
$\mathrm{e}^{-\alpha t}$: evaluate the image function at a complex parameter, then take its
imaginary or real part.

At this point the integral-transform part is complete: the definitions, common pairs and properties
of the Fourier and Laplace transforms, together with the inversion and applications of this chapter.
The next chapter begins the theory of equations of mathematical physics --- the derivation and
solution of the wave, heat and potential equations.

\end{document}
